
\section{Discussion}

\subsection{From Issuing Commands to Stabilizing Language--Action Mappings}

Our findings suggest that verbal robot control involves more than identifying a ``correct'' way to phrase an instruction. Across different strategies, participants sought to make the relationship between what they said and what the robot did sufficiently stable for the interaction to continue. Prior work has shown that conversational partners develop locally shared expressions and that spatial language depends on negotiated or inferred frames of reference. \citep{brennan1996pacts, li2016spatial} Our findings extend this perspective into embodied control. Participants moved between quantitative, qualitative, and relative expressions, introduced locally meaningful units, tested spatial mappings, changed how they controlled continuation and stopping, and reasoned about whether mismatches originated in their own expressions or in robot execution.

Importantly, this uncertainty did not reside in natural language alone. A numerically precise command could still produce an imprecise physical outcome, making linguistic precision distinct from embodied precision. Viewed together, the four dimensions of our taxonomy characterize different aspects of the mapping that users sought to stabilize: magnitude establishes how much, flow control establishes when action continues or stops, spatial reference establishes from whose frame, and error inference helps determine what should be repaired when the mapping fails. Rather than treating these behaviors as isolated command styles, we interpret them as components of a \textit{session-local control convention} through which users make verbal control locally interpretable and manageable.

\subsection{Designing Robots to Participate in Calibration}

Existing work on adaptive and shared human--robot interaction has emphasized that effective collaboration can involve adjustment across both human and robot behavior. \citep{nikolaidis2017mutual, cui2023lilac} In our study, however, much of this adjustment remained on the human side. Participants changed their wording, compensated for observed motion, selected different reference frames, adjusted the granularity of control, and inferred why an interaction had failed.

This shifts the design opportunity from simply building robots that can understand more language toward building robots that can \textit{participate in establishing what language means in the current interaction}. When an expression such as ``a little'' is underspecified, a robot can request clarification rather than silently committing to one interpretation. When users establish a useful local convention, such as a personalized movement unit, the system can retain and revise it. When spatial frames are uncertain, the robot can make its interpretation explicit; when action deviates from the intended outcome, feedback can help users distinguish interpretation errors from execution errors and select an appropriate repair. Recent LLM-enabled robotic systems already demonstrate capabilities such as uncertainty-triggered clarification and learning or retaining user-specific information. \citep{ren2023knowno, wu2023tidybot} These capabilities create an opportunity for verbal control to become less dependent on unilateral human adaptation and more actively negotiated between human and robot.

\subsection{More Capable Robots Make Interaction-Time Grounding More Important, Not Less}

Our findings also help explain why low-level verbal control remains relevant as robots become increasingly capable. Traditional robot interfaces often reduce uncertainty by constraining users to predefined commands, parameters, or interaction rules; much of the mapping between input and action is therefore established by designers before interaction begins. \citep{fong2003questions, marge2010routes} In contrast, language models increasingly allow robots to accept flexible, open-ended, and context-sensitive instructions. \citep{ichter2023saycan, driess2023palme, brohan2023rt2}

This increased flexibility does not eliminate the need for a control protocol. Instead, it shifts part of that protocol \textit{from design time to interaction time}. Expressions such as ``move a little closer,'' ``keep going until I say stop,'' or ``do it like before'' cannot be fully specified independently of the particular user, robot, physical environment, and interaction history. Low-level motion control provides a controlled setting in which this shift becomes especially visible, but the underlying challenge extends beyond navigation. As robots take on richer and more autonomous roles, they will increasingly need to establish, maintain, and repair situated conventions with people rather than merely execute mappings fixed in advance. In this sense, greater language capability does not remove the grounding problem; \textit{it changes where grounding happens and who participates in resolving it}.

